% Irodov, Problems in General Physics, Arihant PDF pp. 21--23,
% printed pp. 12--14. The figures 1.15--1.23 and the pulley constraints
% were checked against the user's original PDF.
\exproblem{71}{12}
电梯轿厢顶上固定一理想滑轮，轻绳两端分别悬挂质量为 $m_1,m_2$ 的重物。
电梯以恒定加速度 $a_0$ 向上启动。求：（a）$m_1$ 相对电梯和相对井道
的加速度；（b）滑轮施加于轿厢顶部的力。
\exsolution
取竖直向下为电梯相对坐标的正方向。在电梯参考系中等效重力加速度
$g_{\rm eff}=g+a_0$，两端相对加速度大小相等且方向相反：
\[
m_1g_{\rm eff}-T=m_1a'_1,\quad
T-m_2g_{\rm eff}=m_2a'_1 .
\]
于是
\[
\boxed{a'_1=\frac{m_1-m_2}{m_1+m_2}(g+a_0)}
\quad\text{（相对电梯，向下为正）}.
\]
井道中的向下加速度需要减去电梯向上的加速度：
\[
\boxed{a_1=a'_1-a_0
 =\frac{(m_1-m_2)g-2m_2a_0}{m_1+m_2}}.
\]
拉力 $T=2m_1m_2(g+a_0)/(m_1+m_2)$，两股绳均竖直向下拉滑轮，
故滑轮对顶部的向下作用力
\[
\boxed{F=2T=\frac{4m_1m_2}{m_1+m_2}(g+a_0)}.
\]
\excheck 若 $m_1=m_2$，两物体相对电梯静止，井道中的加速度为 $-a_0$
（向下取正），与轿厢共同向上运动。底本书末疑将 $(g+a_0)$ 误印为
$(g-a_0)$；由于题面明确写电梯\emph{向上}加速，本稿按动力学推导并记勘误。

\exproblem{72}{12--13}
图1.15中，物块1位于倾角 $\alpha$ 的光滑斜面上，
物块2悬挂于活动滑轮的另一股细绳上；后者质量是前者的
$\eta$ 倍。滑轮、细绳均理想。求物块2的竖直加速度及可能方向。
\exsolution
令 $m_1=m,\ m_2=\eta m$。图中上方绳跨过定滑轮，
它的另一端悬挂活动滑轮；活动滑轮另绕一根一端固定、
另一端悬挂物块2的轻绳。记物块2向下加速度 $a$。
两次不可伸长的绳长关系表明，活动滑轮相应移动 $a/2$，
物块1沿斜面向上的加速度也为 $a/2$。
若第二股绳张力为 $T$，活动滑轮上端绳张力为 $2T$。故
\[
2T-mg\sin\alpha=\frac12ma,\qquad
\eta mg-T=\eta ma.
\]
消去 $T$，得
\[
\boxed{a=\frac{2g(2\eta-\sin\alpha)}{4\eta+1}}.
\]
\excheck 若 $\eta>\tfrac12\sin\alpha$，物块2向下；
若 $\eta<\tfrac12\sin\alpha$，则 $a<0$，物块2向上；
等号时平衡。公式的符号由绳长关系给出，不能预先固定方向。

\exproblem{73}{13}
如图1.16，$m_0$ 在光滑水平面上，绳经固定滑轮拉着下方活动滑轮；
活动滑轮绕另一根细绳，其两端各挂 $m_1,m_2$。
所有滑轮、细绳质量与摩擦可略。求 $m_1$ 的竖直加速度，并讨论方向。
\exsolution
记下方绳张力为 $T$，则连接 $m_0$ 与活动滑轮的张力为 $2T$。
令 $m_1,m_2$ 的竖直向下加速度分别为 $a_1,a_2$，
$m_0$ 向右加速度为 $a_0$。上绳约束给出活动滑轮向下
加速度 $a_p=a_0$；下绳约束给出 $a_1+a_2=2a_p$。由牛顿第二定律
\[
m_1a_1=m_1g-T,\quad m_2a_2=m_2g-T,\quad
m_0a_0=2T.
\]
代入绳长约束：
\[
2g-T\left(\frac1{m_1}+\frac1{m_2}\right)=\frac{4T}{m_0},
\quad
T=\frac{2gm_0m_1m_2}
 {4m_1m_2+m_0(m_1+m_2)}.
\]
因此
\[
\boxed{a_1=
g\,\frac{4m_1m_2+m_0(m_1-m_2)}
 {4m_1m_2+m_0(m_1+m_2)}}.
\]
\excheck $a_1$ 以向下为正；分子为正、零、负分别对应
$m_1$ 向下、静止、向上。若 $m_1=m_2$，二者均向下拉动活动滑轮，
$m_0$ 被拉向右，不应误当成普通阿特伍德机。

\exproblem{74}{13}
图1.17中长为 $\ell$ 的重杆质量为 $M$，
另一端套在绳上并可滑动的小球质量为 $m<M$，
滑动时绳对小球施加未知恒定摩擦力。
从静止释放后，两物体均作匀加速运动；经过 $t$ 秒，
小球由杆下端所处高度变为与杆上端等高。求摩擦力大小。
\exsolution
设绳对球的摩擦力向上为 $f$。
理想绳中小球处上方的张力即 $f$，故杆的向下加速度
$a_M=g-f/M$，小球的向下加速度 $a_m=g-f/m$。
因杆相对小球向下多移动 $\ell$，有
\[
\ell=\frac12(a_M-a_m)t^2
=\frac{ft^2}{2}\left(\frac1m-\frac1M\right).
\]
所以
\[
\boxed{f=\frac{2\ell mM}{(M-m)t^2}}.
\]
\excheck 右端有力的量纲；题设 $M>m$ 保证 $f>0$。
如果摩擦方向与相对滑动方向不符，给定运动便不能实现。

\exproblem{75}{13}
图1.18的复合滑轮中，小球1的质量为杆2质量的
$\eta=1.8$ 倍，杆长 $\ell=1.00\,\mathrm m$。
小球初始位置与杆下端等高，系统由静止释放。
求小球运动至与杆上端等高所需时间（滑轮、细绳理想）。
\exsolution
以小球1向下为正，利用图示两级滑轮的不可伸长关系，
每一级活动滑轮的位移是由两侧绳段变化之和确定的。
设杆的相对小球加速度为 $a_{\rm rel}$。由各质量的牛顿方程、
活动滑轮上 $T_{\rm upper}=2T_{\rm lower}$ 与绳长约束消元，可得
\[
a_{\rm rel}=\frac{3g(2-\eta)}{4+\eta}.
\]
于是从静止开始，两端的相对高度变化满足
$\ell=\frac12a_{\rm rel}t^2$，所以
\[
\boxed{t=\sqrt{\frac{2\ell(4+\eta)}
 {3g(2-\eta)}}\simeq1.40\,\mathrm s}.
\]
\excheck 只有 $\eta<2$ 才有题述的相遇过程；
$\eta\to2^-$ 时所需时间发散。该结果与本版书末简答相符。

\exproblem{76}{13}
图1.19中物块1与2的质量比为 $\eta=m_1/m_2=4.0$。
物块1起初距地面 $h=0.20\,\mathrm m$，
经活动滑轮与物块2相连。由静止释放，
忽略滑轮、绳质量与摩擦。求物块2所能上升的最大高度。
\exsolution
用物块1向下位移 $x$ 作广义坐标，绳长约束给出
物块2向上位移 $2x$。动能为
$K=\frac12(m_1+4m_2)\dot x^2$，重力做功为
$(m_1-2m_2)gx$，所以碰地前物块1向下加速度为
\[
a_1=\frac{m_1-2m_2}{m_1+4m_2}g
=\frac{\eta-2}{\eta+4}g .
\]
物块1下降 $h$ 后落地，物块2已经上升 $2h$，
其瞬时上升速率满足 $v_2^2=(2\dot x)^2=8a_1h$。
物块1停住后，绳子松弛，物块2继续上升
$\Delta H=v_2^2/(2g)=4a_1h/g$，因此
\[
\boxed{H=2h+\Delta H=\frac{6\eta}{\eta+4}h
=0.60\,\mathrm m}.
\]
\excheck 这是以物块2的\emph{初始高度}为基准的上升量；
物块1触地的非弹性冲击会导致上绳松弛，故不能在触地以后
继续套用共同加速度。

\exproblem{77}{14}
如图1.20，质量为 $m$ 的竖直导杆 $A$ 与质量为
$M=\eta m$ 的光滑楔块 $B$ 相接触。导杆仅能沿竖直方向移动，
楔块只能沿水平面移动，所有接触处均无摩擦。
求两者的加速度大小。
\exsolution
用楔块向右位移 $x$ 作广义坐标。斜面上升倾角为 $\alpha$，
图示接触约束使导杆向下位移为 $y=x\tan\alpha$。
于是
\[
K=\frac12M\dot x^2+\frac12m\dot y^2
=\frac12(M+m\tan^2\alpha)\dot x^2,
\]
而导杆重力做正功 $mg\,\mathrm dy=mg\tan\alpha\,\mathrm dx$。
由拉格朗日方程可得
\[
a_B=\ddot x=\frac{mg\tan\alpha}{M+m\tan^2\alpha}
=\boxed{\frac{g}{\tan\alpha+\eta\cot\alpha}},
\]
\[
a_A=\ddot y=a_B\tan\alpha
=\boxed{\frac{g}{1+\eta\cot^2\alpha}}.
\]
\excheck 楔块向右、导杆向下；令 $\eta\to\infty$，两者
加速度都趋零，符合重楔极限。

\exproblem{78}{14}
图1.21中，质量为 $M$ 的楔块在光滑水平面上滑动，
质量为 $m$ 的物块沿楔块的竖直导向面运动，
其间摩擦系数为 $k$；两只滑轮、细绳均理想，
绳的一端按图固定在地面。求物块相对地面的加速度大小。
\exsolution
由图中绳长不变关系可见，物块的竖直位移大小
等于楔块的水平位移大小：取二者相应加速度分量
为 $(-a,-a)$，故物块相对地面的加速度大小为 $\sqrt2a$。
物块与楔块在水平方向同加速度，其竖直导向面提供
水平支持力 $N=ma$；因物块相对楔块向下滑动，
摩擦阻力竖直向上，大小 $kma$。
对运动系统以楔块水平位移为广义坐标，
动能中的有效质量是 $M+2m$，重力沿相应约束做功为 $mg\,\mathrm dx$，
摩擦耗功为 $kma\,\mathrm dx$，于是
\[
(M+2m)a=mg-kma.
\]
消元得
\[
\boxed{a_{\rm block}=\sqrt2 a
=\frac{\sqrt2\,mg}{M+(2+k)m}}.
\]
\excheck $k=0$ 时结果为 $\sqrt2mg/(M+2m)$；
当 $M\to\infty$ 时物块加速度趋零。斜向约束来自绳路，
不可把图形误读为普通固定楔块斜面下滑。

\exproblem{79}{14}
图1.22中，质量相等的物体1与2分别置于大块 $A$
的水平顶面和竖直侧面，用跨过其固定滑轮的绳相连。
两处摩擦系数均为 $k$，求 $A$ 至少须有多大的水平加速度
才能使两物体相对 $A$ 保持静止。
\exsolution
设 $A$ 向右加速度为 $a$，在随 $A$ 加速的坐标系中，
两物体受向左大小为 $ma$ 的惯性力。
物体1顶面法向力 $N_1=mg$，故可用摩擦
$|f_1|\le kmg$；物体2侧面法向力 $N_2=ma$，
故 $|f_2|\le kma$。
沿顶面和竖直侧面的平衡条件为
\[
T-ma+f_1=0,\qquad T-mg+f_2=0 .
\]
存在满足两处摩擦约束的 $T$ 当且仅当
区间 $[ma-kmg,ma+kmg]$ 与
$[mg-kma,mg+kma]$ 相交。
对 $0\le k<1$，最小可行加速度来自
$ma+kmg=mg-kma$，即
\[
\boxed{a_{\min}=\frac{1-k}{1+k}g}.
\]
\excheck 若 $k\ge1$，即使 $a=0$ 静摩擦也足以
在合适张力下保持平衡，故数学上应取
$a_{\min}=\max\{0,(1-k)g/(1+k)\}$。
方向选错会使两处摩擦上限的作用颠倒。

\exproblem{80}{14}
图1.23中，倾角为 $\alpha$ 的棱柱向左作水平加速，
其上放质量 $m$ 的物块，摩擦系数 $k<\cot\alpha$。
求物块相对棱柱保持静止时所允许的最大加速度 $a$。
\exsolution
在棱柱参考系中，向右的惯性力 $ma$ 在斜面切向
向上分量为 $ma\cos\alpha$，而重力切向向下分量为
$mg\sin\alpha$。斜面法向力为
\[
N=m(g\cos\alpha+a\sin\alpha).
\]
当 $a$ 增大到上滑临界时，静摩擦沿斜面向下且大小为 $kN$：
\[
ma\cos\alpha-mg\sin\alpha=kN.
\]
因此
\[
\boxed{a_{\max}
=g\,\frac{\sin\alpha+k\cos\alpha}
{\cos\alpha-k\sin\alpha}
=g\frac{1+k\cot\alpha}{\cot\alpha-k}}.
\]
\excheck 题设 $k<\cot\alpha$ 保证分母为正。
若 $a$ 较小时物块原本有下滑趋势，
还需校核另一个静摩擦下限；本题只要求上限。
